\documentclass[12pt,a4paper]{article}

\usepackage[T2A]{fontenc}
\usepackage[utf8]{inputenc}
\usepackage[russian]{babel}
\usepackage{amsmath,amssymb,amsthm}
\usepackage{enumitem}
\usepackage{listings}
\usepackage{xcolor}
\usepackage{array}
\usepackage{longtable}
\usepackage{hyperref}

\lstdefinestyle{shell}{
    basicstyle=\ttfamily\small,
    breaklines=true,
    frame=single,
    columns=fullflexible,
    keepspaces=true,
    backgroundcolor=\color{gray!5},
    literate=
      {А}{{\CYRA}}1
      {Б}{{\CYRB}}1
      {В}{{\CYRV}}1
      {Г}{{\CYRG}}1
      {Д}{{\CYRD}}1
      {Е}{{\CYRE}}1
      {Ё}{{\CYRYO}}1
      {Ж}{{\CYRZH}}1
      {З}{{\CYRZ}}1
      {И}{{\CYRI}}1
      {Й}{{\CYRISHRT}}1
      {К}{{\CYRK}}1
      {Л}{{\CYRL}}1
      {М}{{\CYRM}}1
      {Н}{{\CYRN}}1
      {О}{{\CYRO}}1
      {П}{{\CYRP}}1
      {Р}{{\CYRR}}1
      {С}{{\CYRS}}1
      {Т}{{\CYRT}}1
      {У}{{\CYRU}}1
      {Ф}{{\CYRF}}1
      {Х}{{\CYRH}}1
      {Ц}{{\CYRC}}1
      {Ч}{{\CYRCH}}1
      {Ш}{{\CYRSH}}1
      {Щ}{{\CYRSHCH}}1
      {Ъ}{{\CYRHRDSN}}1
      {Ы}{{\CYRERY}}1
      {Ь}{{\CYRSFTSN}}1
      {Э}{{\CYREREV}}1
      {Ю}{{\CYRYU}}1
      {Я}{{\CYRYA}}1
      {а}{{\cyra}}1
      {б}{{\cyrb}}1
      {в}{{\cyrv}}1
      {г}{{\cyrg}}1
      {д}{{\cyrd}}1
      {е}{{\cyre}}1
      {ё}{{\cyryo}}1
      {ж}{{\cyrzh}}1
      {з}{{\cyrz}}1
      {и}{{\cyri}}1
      {й}{{\cyrishrt}}1
      {к}{{\cyrk}}1
      {л}{{\cyrl}}1
      {м}{{\cyrm}}1
      {н}{{\cyrn}}1
      {о}{{\cyro}}1
      {п}{{\cyrp}}1
      {р}{{\cyrr}}1
      {с}{{\cyrs}}1
      {т}{{\cyrt}}1
      {у}{{\cyru}}1
      {ф}{{\cyrf}}1
      {х}{{\cyrh}}1
      {ц}{{\cyrc}}1
      {ч}{{\cyrch}}1
      {ш}{{\cyrsh}}1
      {щ}{{\cyrshch}}1
      {ъ}{{\cyrhrdsn}}1
      {ы}{{\cyrery}}1
      {ь}{{\cyrsftsn}}1
      {э}{{\cyrerev}}1
      {ю}{{\cyryu}}1
      {я}{{\cyrya}}1
}

\lstset{style=shell}

\title{Билет №52 \\ \small Инструменты обработки текстовых данных и файловых систем на примере find, sed, grep. (СпБГУ)}
\author{Сериков Владислав Александрович}
\date{16 мая 2026}

\begin{document}

\maketitle
\tableofcontents
\newpage

\section{Инструменты обработки текстовых данных и файловых систем на примере \texttt{find}, \texttt{sed}, \texttt{grep}}

\subsection{Общая идея Unix-инструментов обработки данных}

В Unix-подобных системах значительная часть системного администрирования, анализа данных и разработки выполняется с помощью небольших утилит командной строки. Их философия обычно формулируется так:

\begin{enumerate}
    \item каждая программа делает одну задачу, но делает её хорошо;
    \item программы работают с текстовыми потоками;
    \item программы легко комбинируются через конвейеры;
    \item формат ввода и вывода по возможности остаётся простым и человекочитаемым.
\end{enumerate}

Основными механизмами взаимодействия таких программ являются стандартные потоки:

\begin{itemize}
    \item \texttt{stdin} --- стандартный ввод;
    \item \texttt{stdout} --- стандартный вывод;
    \item \texttt{stderr} --- стандартный поток ошибок.
\end{itemize}

Например:

\begin{verbatim}
grep "error" app.log
\end{verbatim}

ищет строки со словом \texttt{error} в файле \texttt{app.log}, а команда

\begin{verbatim}
grep "error" app.log | sed 's/error/ERROR/g'
\end{verbatim}

сначала выбирает строки, содержащие \texttt{error}, а затем заменяет \texttt{error} на \texttt{ERROR}.

Символ \texttt{|} задаёт конвейер: стандартный вывод одной команды становится стандартным вводом другой.

\subsection{Файловая система и текстовые данные}

\subsubsection{Файловая система}

Файловая система --- это способ организации хранения данных на внешних носителях. В Unix-подобных системах файловая система имеет древовидную структуру. Вершиной дерева является корневой каталог \texttt{/}.

Основные типы объектов файловой системы:

\begin{itemize}
    \item обычный файл;
    \item каталог;
    \item символическая ссылка;
    \item устройство;
    \item именованный канал;
    \item сокет.
\end{itemize}

Путь к файлу бывает:

\begin{itemize}
    \item абсолютным, если начинается с корня, например:
    \begin{verbatim}
    /home/student/project/main.c
    \end{verbatim}

    \item относительным, если задаётся относительно текущего каталога, например:
    \begin{verbatim}
    project/main.c
    \end{verbatim}
\end{itemize}

\subsubsection{Текстовый файл}

Текстовый файл --- это файл, содержимое которого интерпретируется как последовательность строк. Строки обычно отделяются символом перевода строки \texttt{\textbackslash n}.

Многие классические Unix-инструменты ориентированы именно на построчную обработку текстовых файлов.

\subsection{Регулярные выражения}

Команды \texttt{grep} и \texttt{sed} активно используют регулярные выражения.

\subsubsection{Определение}

Регулярное выражение --- это формальный шаблон, описывающий множество строк.

Например, регулярное выражение

\begin{verbatim}
[0-9]+
\end{verbatim}

описывает непустую последовательность десятичных цифр.

\subsubsection{Основные метасимволы}

\begin{center}
\begin{tabular}{|c|p{9cm}|}
\hline
Символ & Значение \\
\hline
\texttt{.} & Любой один символ, кроме перевода строки \\
\hline
\texttt{*} & Ноль или более повторений предыдущего выражения \\
\hline
\texttt{+} & Одно или более повторений предыдущего выражения \\
\hline
\texttt{?} & Ноль или одно повторение предыдущего выражения \\
\hline
\texttt{\^{}} & Начало строки \\
\hline
\texttt{\$} & Конец строки \\
\hline
\texttt{[abc]} & Один символ из множества: \texttt{a}, \texttt{b}, \texttt{c} \\
\hline
\texttt{[\string^abc]} & Один символ, не входящий в указанное множество \\
\hline
\texttt{[0-9]} & Одна цифра \\
\hline
\texttt{(expr)} & Группировка выражения \\
\hline
\texttt{expr1|expr2} & Альтернатива: \texttt{expr1} или \texttt{expr2} \\
\hline
\end{tabular}
\end{center}

Следует различать базовые и расширенные регулярные выражения.

\begin{itemize}
    \item BRE --- basic regular expressions;
    \item ERE --- extended regular expressions.
\end{itemize}

В \texttt{grep} расширенный синтаксис включается опцией \texttt{-E}:

\begin{verbatim}
grep -E "cat|dog" animals.txt
\end{verbatim}

\subsection{\texttt{grep}: поиск строк по шаблону}

\subsubsection{Назначение}

\texttt{grep} --- утилита для поиска строк, соответствующих заданному шаблону. Название исторически связано с командой редактора \texttt{ed}:

\begin{verbatim}
g/re/p
\end{verbatim}

то есть ``global regular expression print''.

Общий вид команды:

\begin{verbatim}
grep [options] pattern [file...]
\end{verbatim}

Если файл не указан, \texttt{grep} читает данные из стандартного ввода.

\subsubsection{Минимальные примеры}

Пусть есть файл \texttt{log.txt}:

\begin{verbatim}
INFO server started
WARNING disk almost full
ERROR cannot open file
INFO request completed
\end{verbatim}

Найти строки с \texttt{ERROR}:

\begin{verbatim}
grep "ERROR" log.txt
\end{verbatim}

Результат:

\begin{verbatim}
ERROR cannot open file
\end{verbatim}

Найти строки, начинающиеся с \texttt{INFO}:

\begin{verbatim}
grep "^INFO" log.txt
\end{verbatim}

Результат:

\begin{verbatim}
INFO server started
INFO request completed
\end{verbatim}

Найти строки, содержащие \texttt{ERROR} или \texttt{WARNING}:

\begin{verbatim}
grep -E "ERROR|WARNING" log.txt
\end{verbatim}

Результат:

\begin{verbatim}
WARNING disk almost full
ERROR cannot open file
\end{verbatim}

\subsubsection{Часто используемые опции \texttt{grep}}

\begin{center}
\begin{tabular}{|c|p{10cm}|}
\hline
Опция & Назначение \\
\hline
\texttt{-i} & Игнорировать регистр символов \\
\hline
\texttt{-v} & Вывести строки, не соответствующие шаблону \\
\hline
\texttt{-n} & Показывать номера строк \\
\hline
\texttt{-c} & Вывести только количество найденных строк \\
\hline
\texttt{-l} & Вывести только имена файлов, где найдено совпадение \\
\hline
\texttt{-r}, \texttt{-R} & Рекурсивный поиск по каталогам \\
\hline
\texttt{-E} & Использовать расширенные регулярные выражения \\
\hline
\texttt{-F} & Интерпретировать шаблон как фиксированную строку, а не регулярное выражение \\
\hline
\texttt{-w} & Искать только целые слова \\
\hline
\texttt{-x} & Совпадение должно занимать всю строку \\
\hline
\texttt{-A n} & Показать \texttt{n} строк после найденной \\
\hline
\texttt{-B n} & Показать \texttt{n} строк до найденной \\
\hline
\texttt{-C n} & Показать \texttt{n} строк контекста до и после найденной \\
\hline
\end{tabular}
\end{center}

\subsubsection{Примеры использования опций}

Поиск без учёта регистра:

\begin{verbatim}
grep -i "error" log.txt
\end{verbatim}

Вывести номера строк:

\begin{verbatim}
grep -n "INFO" log.txt
\end{verbatim}

Результат:

\begin{verbatim}
1:INFO server started
4:INFO request completed
\end{verbatim}

Вывести строки, не содержащие \texttt{INFO}:

\begin{verbatim}
grep -v "INFO" log.txt
\end{verbatim}

Результат:

\begin{verbatim}
WARNING disk almost full
ERROR cannot open file
\end{verbatim}

Подсчитать количество строк с ошибками:

\begin{verbatim}
grep -c "ERROR" log.txt
\end{verbatim}

\subsubsection{Рекурсивный поиск}

Пусть есть каталог проекта:

\begin{verbatim}
project/
    main.py
    utils.py
    README.md
\end{verbatim}

Найти все упоминания функции \texttt{connect}:

\begin{verbatim}
grep -R "connect" project/
\end{verbatim}

Найти только имена файлов, где встречается \texttt{TODO}:

\begin{verbatim}
grep -Rl "TODO" project/
\end{verbatim}

\subsubsection{Алгоритм работы \texttt{grep} на концептуальном уровне}

На высоком уровне \texttt{grep} выполняет следующие шаги:

\begin{enumerate}
    \item Получить шаблон поиска.
    \item Преобразовать регулярное выражение во внутреннее представление.
    \item Открыть входной файл или использовать стандартный ввод.
    \item Читать входной поток построчно.
    \item Для каждой строки проверить, есть ли совпадение с шаблоном.
    \item Если строка удовлетворяет условию поиска, вывести её или выполнить действие, заданное опциями.
    \item После обработки всех строк завершить работу и вернуть код завершения.
\end{enumerate}

\subsubsection{Иллюстрация алгоритма}

Файл:

\begin{verbatim}
1: alpha
2: beta
3: alphabet
4: gamma
\end{verbatim}

Команда:

\begin{verbatim}
grep "alpha" file.txt
\end{verbatim}

Пошаговая обработка:

\begin{center}
\begin{tabular}{|c|c|c|}
\hline
Строка & Проверка шаблона \texttt{alpha} & Действие \\
\hline
\texttt{alpha} & совпадает & вывести \\
\hline
\texttt{beta} & не совпадает & пропустить \\
\hline
\texttt{alphabet} & совпадает & вывести \\
\hline
\texttt{gamma} & не совпадает & пропустить \\
\hline
\end{tabular}
\end{center}

Результат:

\begin{verbatim}
alpha
alphabet
\end{verbatim}

\subsubsection{Коды завершения \texttt{grep}}

Обычно \texttt{grep} возвращает:

\begin{itemize}
    \item \texttt{0}, если совпадение найдено;
    \item \texttt{1}, если совпадение не найдено;
    \item \texttt{2}, если произошла ошибка.
\end{itemize}

Это удобно использовать в сценариях командной оболочки:

\begin{verbatim}
if grep -q "ERROR" log.txt; then
    echo "Errors found"
else
    echo "No errors"
fi
\end{verbatim}

Опция \texttt{-q} подавляет вывод и используется только для проверки факта совпадения.

\subsection{\texttt{sed}: потоковый редактор}

\subsubsection{Назначение}

\texttt{sed} --- потоковый редактор. Он читает входной поток, применяет к нему заданный сценарий редактирования и выводит результат.

Общий вид:

\begin{verbatim}
sed [options] 'script' [file...]
\end{verbatim}

Ключевая особенность \texttt{sed}: он обрабатывает данные последовательно, строка за строкой, не требуя загрузки всего файла в память.

\subsubsection{Основной цикл работы \texttt{sed}}

Типичный цикл работы \texttt{sed} можно описать так:

\begin{enumerate}
    \item Прочитать очередную строку из входного потока в рабочую область, называемую pattern space.
    \item Применить к ней команды сценария.
    \item Если автоматическая печать не отключена, вывести содержимое pattern space.
    \item Очистить pattern space.
    \item Перейти к следующей строке.
\end{enumerate}

\subsubsection{Иллюстрация цикла}

Файл \texttt{names.txt}:

\begin{verbatim}
Ann
Bob
Alice
\end{verbatim}

Команда:

\begin{verbatim}
sed 's/A/a/' names.txt
\end{verbatim}

Пошагово:

\begin{center}
\begin{tabular}{|c|c|c|}
\hline
Входная строка & После команды \texttt{s/A/a/} & Вывод \\
\hline
\texttt{Ann} & \texttt{ann} & \texttt{ann} \\
\hline
\texttt{Bob} & \texttt{Bob} & \texttt{Bob} \\
\hline
\texttt{Alice} & \texttt{alice} & \texttt{alice} \\
\hline
\end{tabular}
\end{center}

Результат:

\begin{verbatim}
ann
Bob
alice
\end{verbatim}

\subsubsection{Команда подстановки}

Наиболее часто используется команда \texttt{s}:

\begin{verbatim}
s/pattern/replacement/flags
\end{verbatim}

где:

\begin{itemize}
    \item \texttt{pattern} --- регулярное выражение;
    \item \texttt{replacement} --- строка замены;
    \item \texttt{flags} --- дополнительные флаги.
\end{itemize}

Пример:

\begin{verbatim}
echo "one two one" | sed 's/one/ONE/'
\end{verbatim}

Результат:

\begin{verbatim}
ONE two one
\end{verbatim}

По умолчанию заменяется только первое вхождение в строке. Чтобы заменить все вхождения, используется флаг \texttt{g}:

\begin{verbatim}
echo "one two one" | sed 's/one/ONE/g'
\end{verbatim}

Результат:

\begin{verbatim}
ONE two ONE
\end{verbatim}

\subsubsection{Адресация строк в \texttt{sed}}

Команды \texttt{sed} можно применять не ко всем строкам, а только к выбранным.

Примеры адресов:

\begin{center}
\begin{tabular}{|c|p{9cm}|}
\hline
Адрес & Значение \\
\hline
\texttt{3} & Только третья строка \\
\hline
\texttt{1,5} & Строки с первой по пятую \\
\hline
\texttt{\$} & Последняя строка \\
\hline
\texttt{/regex/} & Строки, соответствующие регулярному выражению \\
\hline
\end{tabular}
\end{center}

Пример: заменить \texttt{error} на \texttt{ERROR} только в строках, содержащих \texttt{critical}:

\begin{verbatim}
sed '/critical/s/error/ERROR/g' log.txt
\end{verbatim}

Удалить вторую строку:

\begin{verbatim}
sed '2d' file.txt
\end{verbatim}

Вывести только строки с 10-й по 20-ю:

\begin{verbatim}
sed -n '10,20p' file.txt
\end{verbatim}

Здесь опция \texttt{-n} отключает автоматическую печать, а команда \texttt{p} явно печатает выбранные строки.

\subsubsection{Основные команды \texttt{sed}}

\begin{center}
\begin{tabular}{|c|p{10cm}|}
\hline
Команда & Назначение \\
\hline
\texttt{p} & Напечатать строку \\
\hline
\texttt{d} & Удалить строку из вывода \\
\hline
\texttt{s} & Выполнить замену \\
\hline
\texttt{a} & Добавить строку после текущей \\
\hline
\texttt{i} & Вставить строку перед текущей \\
\hline
\texttt{c} & Заменить выбранную строку целиком \\
\hline
\texttt{q} & Завершить обработку \\
\hline
\texttt{y} & Транслитерировать символы \\
\hline
\end{tabular}
\end{center}

\subsubsection{Примеры редактирования}

Удалить пустые строки:

\begin{verbatim}
sed '/^$/d' file.txt
\end{verbatim}

Удалить комментарии, начинающиеся с \texttt{\#}:

\begin{verbatim}
sed '/^#/d' config.conf
\end{verbatim}

Заменить несколько пробелов одним пробелом:

\begin{verbatim}
sed -E 's/[ ]+/ /g' file.txt
\end{verbatim}

Добавить строку после первой:

\begin{verbatim}
sed '1a inserted line' file.txt
\end{verbatim}

Вставить строку перед первой:

\begin{verbatim}
sed '1i header line' file.txt
\end{verbatim}

\subsubsection{Редактирование файла на месте}

По умолчанию \texttt{sed} не изменяет исходный файл, а выводит результат в стандартный вывод.

Например:

\begin{verbatim}
sed 's/foo/bar/g' file.txt
\end{verbatim}

не изменяет \texttt{file.txt}. Чтобы записать результат в файл, можно использовать перенаправление:

\begin{verbatim}
sed 's/foo/bar/g' file.txt > new_file.txt
\end{verbatim}

В GNU \texttt{sed} есть опция \texttt{-i}, которая изменяет файл на месте:

\begin{verbatim}
sed -i 's/foo/bar/g' file.txt
\end{verbatim}

Безопаснее делать резервную копию:

\begin{verbatim}
sed -i.bak 's/foo/bar/g' file.txt
\end{verbatim}

После этого появится файл \texttt{file.txt.bak}.

\subsubsection{Группы и обратные ссылки}

Пусть есть файл:

\begin{verbatim}
Ivan Petrov
Anna Sidorova
\end{verbatim}

Нужно поменять имя и фамилию местами:

\begin{verbatim}
sed -E 's/([A-Za-z]+) ([A-Za-z]+)/\2, \1/' names.txt
\end{verbatim}

Результат:

\begin{verbatim}
Petrov, Ivan
Sidorova, Anna
\end{verbatim}

Здесь:

\begin{itemize}
    \item \texttt{([A-Za-z]+)} --- первая группа;
    \item пробел разделяет имя и фамилию;
    \item \texttt{([A-Za-z]+)} --- вторая группа;
    \item \texttt{\textbackslash 1} --- содержимое первой группы;
    \item \texttt{\textbackslash 2} --- содержимое второй группы.
\end{itemize}

\subsection{\texttt{find}: поиск объектов файловой системы}

\subsubsection{Назначение}

\texttt{find} --- утилита для обхода дерева каталогов и поиска файлов по заданным условиям.

Общий вид:

\begin{verbatim}
find [path...] [expression]
\end{verbatim}

Примеры:

\begin{verbatim}
find .
find /home/student -name "*.txt"
find . -type f -size +10M
\end{verbatim}

Если путь не указан, обычно используется текущий каталог. Точка \texttt{.} явно обозначает текущий каталог.

\subsubsection{Основные критерии поиска}

\begin{center}
\begin{tabular}{|c|p{10cm}|}
\hline
Критерий & Назначение \\
\hline
\texttt{-name pattern} & Поиск по имени файла \\
\hline
\texttt{-iname pattern} & Поиск по имени без учёта регистра \\
\hline
\texttt{-type f} & Обычные файлы \\
\hline
\texttt{-type d} & Каталоги \\
\hline
\texttt{-type l} & Символические ссылки \\
\hline
\texttt{-size n} & Поиск по размеру \\
\hline
\texttt{-mtime n} & Поиск по времени изменения в днях \\
\hline
\texttt{-user name} & Файлы указанного пользователя \\
\hline
\texttt{-group name} & Файлы указанной группы \\
\hline
\texttt{-perm mode} & Поиск по правам доступа \\
\hline
\end{tabular}
\end{center}

\subsubsection{Минимальные примеры}

Найти все файлы с расширением \texttt{.txt} в текущем каталоге и подкаталогах:

\begin{verbatim}
find . -name "*.txt"
\end{verbatim}

Найти все обычные файлы:

\begin{verbatim}
find . -type f
\end{verbatim}

Найти все каталоги:

\begin{verbatim}
find . -type d
\end{verbatim}

Найти файлы больше 100 мегабайт:

\begin{verbatim}
find . -type f -size +100M
\end{verbatim}

Найти файлы, изменённые за последние 7 дней:

\begin{verbatim}
find . -type f -mtime -7
\end{verbatim}

Найти файлы, изменённые более 30 дней назад:

\begin{verbatim}
find . -type f -mtime +30
\end{verbatim}

\subsubsection{Действия \texttt{find}}

\texttt{find} не только ищет файлы, но и выполняет действия над найденными объектами.

\begin{center}
\begin{tabular}{|c|p{10cm}|}
\hline
Действие & Назначение \\
\hline
\texttt{-print} & Напечатать путь \\
\hline
\texttt{-delete} & Удалить найденные объекты \\
\hline
\texttt{-exec command \{\} \textbackslash;} & Выполнить команду для каждого найденного объекта \\
\hline
\texttt{-exec command \{\} +} & Выполнить команду с группой найденных объектов \\
\hline
\end{tabular}
\end{center}

Пример: удалить временные файлы:

\begin{verbatim}
find . -name "*.tmp" -type f -delete
\end{verbatim}

Пример: показать подробную информацию о найденных файлах:

\begin{verbatim}
find . -name "*.txt" -exec ls -l {} \;
\end{verbatim}

Пример: передать несколько файлов одной команде:

\begin{verbatim}
find . -name "*.txt" -exec wc -l {} +
\end{verbatim}

Разница:

\begin{itemize}
    \item \texttt{\{\} \textbackslash;} --- команда запускается отдельно для каждого файла;
    \item \texttt{\{\} +} --- команда запускается для группы файлов, что обычно эффективнее.
\end{itemize}

\subsubsection{Логические операции в \texttt{find}}

Выражения \texttt{find} можно комбинировать логическими операторами.

\begin{center}
\begin{tabular}{|c|p{9cm}|}
\hline
Оператор & Значение \\
\hline
\texttt{-a} & Логическое И, часто подразумевается неявно \\
\hline
\texttt{-o} & Логическое ИЛИ \\
\hline
\texttt{!} & Логическое НЕ \\
\hline
\texttt{( expr )} & Группировка условий \\
\hline
\end{tabular}
\end{center}

Пример: найти \texttt{.c} и \texttt{.h} файлы:

\begin{verbatim}
find . \( -name "*.c" -o -name "*.h" \) -type f
\end{verbatim}

Скобки экранируются, потому что иначе их может интерпретировать командная оболочка.

Найти все файлы, кроме \texttt{.log}:

\begin{verbatim}
find . -type f ! -name "*.log"
\end{verbatim}

\subsubsection{Алгоритм обхода дерева каталогов}

Утилита \texttt{find} концептуально выполняет обход дерева файловой системы.

Один из возможных вариантов обхода --- поиск в глубину.

\paragraph{Алгоритм обхода в глубину для \texttt{find}.}

\begin{enumerate}
    \item Поместить начальный каталог в стек.
    \item Пока стек не пуст:
    \begin{enumerate}
        \item извлечь очередной путь;
        \item получить информацию об объекте файловой системы;
        \item проверить объект на соответствие выражению поиска;
        \item если объект удовлетворяет условиям, выполнить заданное действие;
        \item если объект является каталогом, добавить его содержимое в стек.
    \end{enumerate}
    \item Завершить работу после обработки всех доступных объектов.
\end{enumerate}

\subsubsection{Иллюстрация обхода}

Пусть дерево каталогов имеет вид:

\begin{verbatim}
project/
    README.md
    src/
        main.c
        util.c
    tests/
        test_main.c
\end{verbatim}

Команда:

\begin{verbatim}
find project -name "*.c"
\end{verbatim}

Возможная последовательность обработки:

\begin{enumerate}
    \item проверить \texttt{project};
    \item проверить \texttt{project/README.md};
    \item перейти в \texttt{project/src};
    \item проверить \texttt{project/src/main.c};
    \item проверить \texttt{project/src/util.c};
    \item перейти в \texttt{project/tests};
    \item проверить \texttt{project/tests/test\_main.c}.
\end{enumerate}

Результат:

\begin{verbatim}
project/src/main.c
project/src/util.c
project/tests/test_main.c
\end{verbatim}

\subsubsection{Ограничение глубины поиска}

Иногда нужно искать не во всём дереве каталогов, а только на определённой глубине.

Найти файлы только в текущем каталоге, без подкаталогов:

\begin{verbatim}
find . -maxdepth 1 -type f
\end{verbatim}

Искать начиная с глубины 2:

\begin{verbatim}
find . -mindepth 2 -type f
\end{verbatim}

\subsubsection{Работа с датами}

Найти файлы, изменённые ровно 1 день назад:

\begin{verbatim}
find . -mtime 1
\end{verbatim}

Найти файлы, изменённые менее суток назад:

\begin{verbatim}
find . -mtime -1
\end{verbatim}

Найти файлы, изменённые более 10 дней назад:

\begin{verbatim}
find . -mtime +10
\end{verbatim}

Полезные критерии времени:

\begin{itemize}
    \item \texttt{-mtime} --- время последнего изменения содержимого файла;
    \item \texttt{-atime} --- время последнего доступа;
    \item \texttt{-ctime} --- время последнего изменения метаданных.
\end{itemize}

\subsection{Совместное использование \texttt{find}, \texttt{grep} и \texttt{sed}}

\subsubsection{Поиск файлов и поиск текста внутри них}

Найти все файлы \texttt{.py}, содержащие строку \texttt{import requests}:

\begin{verbatim}
find . -name "*.py" -type f -exec grep -n "import requests" {} +
\end{verbatim}

Здесь:

\begin{itemize}
    \item \texttt{find .} ищет в текущем каталоге;
    \item \texttt{-name "*.py"} выбирает Python-файлы;
    \item \texttt{-type f} оставляет только обычные файлы;
    \item \texttt{-exec grep -n ... {} +} запускает \texttt{grep} для найденных файлов.
\end{itemize}

\subsubsection{Массовая замена в найденных файлах}

Заменить \texttt{localhost} на \texttt{127.0.0.1} во всех \texttt{.conf}-файлах:

\begin{verbatim}
find . -name "*.conf" -type f -exec sed -i.bak 's/localhost/127.0.0.1/g' {} +
\end{verbatim}

При этом создаются резервные копии с расширением \texttt{.bak}.

\subsubsection{Конвейерная обработка}

Вывести строки с ошибками из всех логов, заменив \texttt{ERROR} на \texttt{ERR}:

\begin{verbatim}
find . -name "*.log" -type f -exec grep "ERROR" {} + | sed 's/ERROR/ERR/g'
\end{verbatim}

\subsection{Проблема пробелов и специальных символов в именах файлов}

Имена файлов могут содержать пробелы, табуляции и даже символы перевода строки. Поэтому конструкции вида:

\begin{verbatim}
for f in $(find . -name "*.txt"); do
    echo "$f"
done
\end{verbatim}

опасны: результат \texttt{find} разбивается оболочкой по пробельным символам.

Более безопасный способ:

\begin{verbatim}
find . -name "*.txt" -type f -print0 |
while IFS= read -r -d '' file; do
    echo "$file"
done
\end{verbatim}

Здесь:

\begin{itemize}
    \item \texttt{-print0} отделяет имена файлов нулевым байтом;
    \item \texttt{read -d ''} читает до нулевого байта;
    \item \texttt{IFS=} отключает разбиение по пробелам;
    \item \texttt{-r} запрещает специальную обработку обратной косой черты.
\end{itemize}

Также часто используется связка с \texttt{xargs -0}:

\begin{verbatim}
find . -name "*.txt" -type f -print0 | xargs -0 grep "pattern"
\end{verbatim}

\subsection{Сравнение \texttt{find}, \texttt{grep}, \texttt{sed}}

\begin{center}
\begin{tabular}{|c|p{4cm}|p{6cm}|}
\hline
Утилита & Основная задача & Типичный пример \\
\hline
\texttt{find} & Поиск файлов и каталогов по свойствам & Найти все \texttt{.log}-файлы старше 30 дней \\
\hline
\texttt{grep} & Поиск строк в текстовом потоке & Найти строки с \texttt{ERROR} \\
\hline
\texttt{sed} & Потоковое редактирование текста & Заменить \texttt{foo} на \texttt{bar} \\
\hline
\end{tabular}
\end{center}

\subsection{Типовые практические задачи}

\subsubsection{Найти большие файлы}

\begin{verbatim}
find . -type f -size +1G
\end{verbatim}

\subsubsection{Найти и удалить пустые каталоги}

\begin{verbatim}
find . -type d -empty -delete
\end{verbatim}

\subsubsection{Найти строки с IPv4-адресами}

\begin{verbatim}
grep -E "([0-9]{1,3}\.){3}[0-9]{1,3}" file.txt
\end{verbatim}

Это выражение синтаксически находит последовательности вида \texttt{192.168.0.1}, но не проверяет строго, что каждое число лежит в диапазоне от 0 до 255.

\subsubsection{Удалить завершающие пробелы}

\begin{verbatim}
sed 's/[[:space:]]*$//' file.txt
\end{verbatim}

\subsubsection{Переименовать расширение через \texttt{find} и shell}

Например, заменить расширение \texttt{.txt} на \texttt{.bak}:

\begin{verbatim}
find . -type f -name "*.txt" -exec sh -c '
for file do
    mv "$file" "${file%.txt}.bak"
done
' sh {} +
\end{verbatim}

Здесь конструкция \texttt{\$\{file\%.txt\}} удаляет суффикс \texttt{.txt} из значения переменной \texttt{file}.

\subsection{Теоретические замечания о регулярных выражениях}

\subsubsection{Регулярные языки}

\textbf{Определение.}
Алфавит --- конечное множество символов. Слово над алфавитом --- конечная последовательность символов этого алфавита. Язык --- множество слов над заданным алфавитом.

\textbf{Определение.}
Класс регулярных языков над алфавитом $\Sigma$ определяется индуктивно:

\begin{enumerate}
    \item $\varnothing$ является регулярным языком;
    \item $\{\varepsilon\}$ является регулярным языком, где $\varepsilon$ --- пустое слово;
    \item для каждого $a \in \Sigma$ язык $\{a\}$ является регулярным;
    \item если $L_1$ и $L_2$ регулярны, то $L_1 \cup L_2$ регулярный;
    \item если $L_1$ и $L_2$ регулярны, то $L_1 L_2$ регулярный, где
    \[
        L_1 L_2 = \{xy \mid x \in L_1, y \in L_2\};
    \]
    \item если $L$ регулярен, то $L^*$ регулярный, где
    \[
        L^* = \bigcup_{n=0}^{\infty} L^n.
    \]
\end{enumerate}

Регулярные выражения являются синтаксическим способом задания регулярных языков.

\subsubsection{Теорема Клини}

\textbf{Теорема Клини.}
Язык является регулярным тогда и только тогда, когда он распознаётся некоторым конечным автоматом.

\textbf{Доказательство.}

Докажем обе стороны.

\paragraph{1. Если язык задаётся регулярным выражением, то он распознаётся конечным автоматом.}

Доказательство проводится структурной индукцией по построению регулярного выражения.

\begin{enumerate}
    \item Для языка $\varnothing$ строится автомат без принимающих состояний.
    \item Для языка $\{\varepsilon\}$ строится автомат, начальное состояние которого является принимающим.
    \item Для языка $\{a\}$ строится автомат с переходом по символу $a$ из начального состояния в принимающее.
\end{enumerate}

Теперь покажем замкнутость относительно операций регулярных выражений.

Пусть языки $L_1$ и $L_2$ распознаются конечными автоматами $A_1$ и $A_2$.

\begin{itemize}
    \item Для объединения $L_1 \cup L_2$ строится новый автомат с новым начальным состоянием, из которого есть $\varepsilon$-переходы в начальные состояния автоматов $A_1$ и $A_2$. Принимающими состояниями являются принимающие состояния обоих автоматов.
    \item Для конкатенации $L_1L_2$ добавляются $\varepsilon$-переходы из каждого принимающего состояния автомата $A_1$ в начальное состояние автомата $A_2$. Начальным состоянием становится начальное состояние $A_1$, принимающими --- принимающие состояния $A_2$.
    \item Для итерации $L_1^*$ добавляется новое начальное принимающее состояние. Из него проводится $\varepsilon$-переход в начальное состояние $A_1$, а из принимающих состояний $A_1$ проводятся $\varepsilon$-переходы обратно в начальное состояние $A_1$ и в новое принимающее состояние.
\end{itemize}

Таким образом, для любого регулярного выражения можно построить конечный автомат, распознающий тот же язык.

\paragraph{2. Если язык распознаётся конечным автоматом, то он задаётся регулярным выражением.}

Пусть дан конечный автомат
\[
    A = (Q, \Sigma, \delta, q_0, F),
\]
где $Q$ --- конечное множество состояний, $\Sigma$ --- алфавит, $\delta$ --- функция переходов, $q_0$ --- начальное состояние, $F$ --- множество принимающих состояний.

Идея доказательства состоит в последовательном исключении состояний автомата с сохранением языка.

Рассмотрим обобщённый конечный автомат, в котором переходы помечаются регулярными выражениями, а не отдельными символами. Добавим новое начальное состояние $s$ и новое финальное состояние $t$. Проведём $\varepsilon$-переход из $s$ в старое начальное состояние $q_0$, а из каждого старого принимающего состояния проведём $\varepsilon$-переход в $t$.

Далее будем удалять промежуточные состояния. Пусть удаляется состояние $r$. Для каждой пары состояний $i$ и $j$, отличных от $r$, обновим метку перехода из $i$ в $j$ по формуле:
\[
    R_{ij}' = R_{ij} \cup R_{ir}(R_{rr})^*R_{rj}.
\]

Смысл формулы следующий:

\begin{itemize}
    \item $R_{ij}$ описывает все пути из $i$ в $j$, не проходящие через $r$;
    \item $R_{ir}(R_{rr})^*R_{rj}$ описывает пути, которые идут из $i$ в $r$, затем любое число раз проходят цикл из $r$ в $r$, а затем идут из $r$ в $j$.
\end{itemize}

После удаления всех промежуточных состояний останутся только $s$ и $t$, а метка перехода из $s$ в $t$ будет регулярным выражением, описывающим ровно язык автомата.

Следовательно, всякий язык, распознаваемый конечным автоматом, задаётся регулярным выражением.

Обе импликации доказаны, значит теорема доказана.

\subsection{Практические рекомендации}

\begin{enumerate}
    \item Для простого поиска фиксированной строки использовать:
    \begin{verbatim}
    grep -F "text" file
    \end{verbatim}

    \item Для поиска по регулярному выражению использовать:
    \begin{verbatim}
    grep -E "regex" file
    \end{verbatim}

    \item Для безопасного изменения файлов через \texttt{sed} сначала делать резервную копию:
    \begin{verbatim}
    sed -i.bak 's/old/new/g' file
    \end{verbatim}

    \item Для поиска файлов по имени использовать \texttt{find}:
    \begin{verbatim}
    find . -name "*.txt"
    \end{verbatim}

    \item Для обработки имён файлов с пробелами использовать \texttt{-print0} и \texttt{xargs -0}:
    \begin{verbatim}
    find . -type f -print0 | xargs -0 grep "pattern"
    \end{verbatim}

    \item Перед массовым удалением сначала проверять список:
    \begin{verbatim}
    find . -name "*.tmp" -type f -print
    \end{verbatim}

    и только затем выполнять:
    \begin{verbatim}
    find . -name "*.tmp" -type f -delete
    \end{verbatim}
\end{enumerate}

\subsection{Краткое резюме}

\begin{itemize}
    \item \texttt{grep} используется для поиска строк по шаблону.
    \item \texttt{sed} используется для потокового редактирования текста.
    \item \texttt{find} используется для поиска объектов файловой системы по имени, типу, размеру, времени изменения, правам доступа и другим признакам.
    \item Эти утилиты особенно эффективны в комбинации друг с другом через конвейеры и \texttt{-exec}.
    \item Регулярные выражения дают формальный язык описания множеств строк и лежат в основе многих операций поиска и замены.
    \item При работе с именами файлов следует учитывать пробелы и специальные символы; безопасными являются подходы с \texttt{-print0} и \texttt{xargs -0}.
\end{itemize}

\end{document}
